% ============================================================
%  PRUEBAS INTEGRALES  (CP-I001 – CP-I008)
% ============================================================
 
\section{Pruebas Integrales}

\subsection{CP-I001: Flujo completo --- crear producto y venderlo}
\def\cpPasos{%
  \item Crear categoría ``Bisutería'' (código BIS-001).
  \item Crear producto ``Aretes de plata'' asociado a la categoría, precio \CRC{}2500.
  \item Intentar vender sin stock; observar respuesta.
  \item Registrar entrada de 10 unidades en Inventario.
  \item Crear venta de 2 unidades de ``Aretes de plata''.
  \item Verificar stock final en Inventario.}
\def\cpResultadoE{%
  \item Stock final = 8 (10 entrada $-$ 2 venta).
  \item Sistema bloquea venta sin stock.}
\def\cpResultadoO{Al intentar vender con stock 0, el sistema mostró error ``stock
  insuficiente para el movimiento''. Luego de registrar 10 unidades, la venta de 2
  unidades se completó correctamente. Stock final = 8. Flujo completo funcionó.}
\casoprueba
  {CP-I001}{2026-04-03}{Categorías \arr{} Productos \arr{} Inventario \arr{} Ventas}
  {Flujo completo de creación de producto hasta venta}
  {Integral}{\estadoAprobado}
  {\item La aplicación está ejecutándose.
   \item Usuario autenticado con \texttt{test@inventario.local / test123}.}
  {\item Categoría: Bisutería, código BIS-001.
   \item Producto: Aretes de plata, código ARE-001, precio \CRC{}2500.
   \item Entrada de stock: 10 unidades. Venta: 2 unidades.}
  {Angel Segura Méndez}

\subsection{CP-I002: Flujo completo --- crear usuario, asignar permisos y verificar acceso}
\def\cpPasos{%
  \item Navegar a ``Usuarios'' \arr{} crear nuevo usuario con los datos indicados.
  \item Asignar solo permiso ``Administración''.
  \item Guardar usuario.
  \item Cerrar sesión.
  \item Iniciar sesión con \texttt{admin@test.com / test123}.
  \item Verificar módulos disponibles en el menú.}
\def\cpResultadoE{%
  \item Menú muestra solo módulos administrativos.
  \item No muestra Ventas, Estadísticas ni Reportes.}
\def\cpResultadoO{El menú mostró exactamente 7 módulos: Inicio, Productos, Categorías,
  Emprendimientos, Emprendedoras, Actividades y Asociación. Los módulos de ventas no
  aparecieron. El sistema aplicó correctamente los permisos asignados.}
\casoprueba
  {CP-I002}{2026-04-03}{Usuarios \arr{} Permisos \arr{} Autenticación \arr{} Navegación}
  {Flujo completo de gestión de usuario con permisos}
  {Integral}{\estadoAprobado}
  {\item La aplicación está ejecutándose.
   \item Usuario administrador autenticado.}
  {\item Usuario: Ana Prueba, correo \texttt{admin@test.com}, contraseña \texttt{test123}.
   \item Permiso asignado: solo Administración.}
  {Angel Segura Méndez}

\subsection{CP-I003: Flujo completo --- generación y exportación de reporte PDF}
\def\cpPasos{%
  \item Navegar a ``Reportes''.
  \item Seleccionar rango de fechas que incluya hoy.
  \item Generar el reporte.
  \item Exportar a PDF.}
\def\cpResultadoE{%
  \item El PDF se genera correctamente con los datos de ventas del período.}
\def\cpResultadoO{El reporte se generó y el PDF se exportó correctamente.}
\casoprueba
  {CP-I003}{2026-04-03}{Ventas \arr{} Reportes \arr{} PDF}
  {Generación y exportación de reporte de ventas en PDF}
  {Integral}{\estadoAprobado}
  {\item La aplicación está ejecutándose.
   \item Existen ventas registradas en el sistema.}
  {\item Rango de fechas: incluye 2026-04-03.}
  {Angel Segura Méndez}

\subsection{CP-I004: Flujo completo --- crear actividad/taller}
\def\cpPasos{%
  \item Navegar a ``Actividades'' \arr{} nueva actividad.
  \item Ingresar título ``Taller de tejido'', tipo y fecha inicio.
  \item Guardar.
  \item Verificar que aparece en la lista de actividades.}
\def\cpResultadoE{%
  \item La actividad ``Taller de tejido'' aparece en la lista.}
\def\cpResultadoO{La actividad se creó y apareció correctamente en la lista de actividades.}
\casoprueba
  {CP-I004}{2026-04-03}{Actividades}
  {Crear actividad y verificar que aparece en la lista}
  {Integral}{\estadoAprobado}
  {\item La aplicación está ejecutándose.
   \item Usuario autenticado.}
  {\item Título: Taller de tejido.
   \item Tipo: Taller. Fecha inicio: 2026-04-03.}
  {Angel Segura Méndez}

\subsection{CP-I005: Editar producto y verificar protección de integridad al eliminar}
\def\cpPasos{%
  \item Navegar a ``Productos'' \arr{} editar ``Aretes de plata''.
  \item Cambiar precio a \CRC{}3000.
  \item Guardar y verificar que el precio se actualizó en la lista.
  \item Intentar eliminar ``Aretes de plata''.}
\def\cpResultadoE{%
  \item La edición funciona correctamente.
  \item La eliminación falla con mensaje por registros asociados.}
\def\cpResultadoO{El precio se actualizó correctamente a \CRC{}3000. Al intentar
  eliminar, el sistema mostró: ``no se puede eliminar porque tiene registros asociados''.
  La integridad referencial funciona correctamente.}
\casoprueba
  {CP-I005}{2026-04-03}{Productos}
  {Editar producto y verificar protección de integridad al eliminar}
  {Integral}{\estadoAprobado}
  {\item Existe el producto ``Aretes de plata'' con ventas asociadas.}
  {\item Producto: Aretes de plata. Nuevo precio: \CRC{}3000.}
  {Angel Segura Méndez}

\subsection{CP-I006: Edición de información de la Asociación}
\def\cpPasos{%
  \item Navegar a ``Asociación''.
  \item Verificar que hay información cargada.
  \item Editar el campo de Misión.
  \item Guardar los cambios.
  \item Verificar que el cambio persiste.}
\def\cpResultadoE{%
  \item La información se actualiza y persiste correctamente.}
\def\cpResultadoO{Había información cargada en el módulo. El cambio en la misión
  se guardó correctamente.}
\casoprueba
  {CP-I006}{2026-04-03}{Asociación}
  {Editar y guardar información de la asociación}
  {Integral}{\estadoAprobado}
  {\item La aplicación está ejecutándose.
   \item Existe información de la asociación cargada.}
  {\item Campo editado: Misión (texto modificado).}
  {Angel Segura Méndez}

\subsection{CP-I007: Búsqueda y filtro de productos}
\def\cpPasos{%
  \item Navegar a ``Productos''.
  \item Buscar por texto ``pulsera''.
  \item Verificar resultados.
  \item Filtrar por categoría ``Artesanías''.
  \item Verificar resultados del filtro.}
\def\cpResultadoE{%
  \item La búsqueda y el filtro muestran solo los productos que coinciden.}
\def\cpResultadoO{Tanto la búsqueda por nombre ``pulsera'' como el filtro por
  categoría ``Artesanías'' funcionaron correctamente, mostrando solo los productos
  correspondientes.}
\casoprueba
  {CP-I007}{2026-04-03}{Productos}
  {Búsqueda por nombre y filtro por categoría funcionan correctamente}
  {Integral}{\estadoAprobado}
  {\item Existen productos en el sistema, incluyendo ``Pulsera Artesanal''.
   \item Existe la categoría ``Artesanías'' con productos asociados.}
  {\item Búsqueda: ``pulsera''. Filtro por categoría: Artesanías.}
  {Angel Segura Méndez}

\subsection{CP-I008: Backup de la base de datos}
\def\cpPasos{%
  \item Localizar la opción de Backup en el menú.
  \item Ejecutar el backup.
  \item Seleccionar carpeta de destino.
  \item Verificar que el archivo se generó.}
\def\cpResultadoE{%
  \item El sistema genera el archivo de respaldo SQLite en la carpeta seleccionada.}
\def\cpResultadoO{El sistema solicitó la carpeta de destino y generó el respaldo de
  la base de datos SQLite correctamente.}
\casoprueba
  {CP-I008}{2026-04-03}{Backup}
  {Generación de respaldo de la base de datos SQLite}
  {Integral}{\estadoAprobado}
  {\item La aplicación está ejecutándose.
   \item Usuario autenticado con permisos suficientes.}
  {\item Carpeta de destino seleccionada por el usuario.}
  {Angel Segura Méndez}
